\documentclass[../../main.tex]{subfiles}
\begin{document}

\begin{center}
  {\fontsize{48pt}{58pt}\selectfont \textbf{1998}\normalsize\textbf{年度}}
\end{center}

\begin{table}[htb]
  \centering
  \caption{出題テーマと難易度}
  \begin{tabular}{|c|p{9cm}|c|}
    \hline
    \textbf{問題番号} & \textbf{テーマ} & \textbf{難易度} \\
    \hline
    第1問 & 内接円をもつ直角三角形の面積 & \\
    \hline
    第2問 & 二の冪による整除性 & \\
    \hline
    第3問 & 四面体の辺上にある平行四辺形 & \\
    \hline
    第4問 & 放物線内の格子点数と面積の比 & \\
    \hline
    第5問 & 色と番号をもつ玉の得点分布 & \\
    \hline
    第6問 & 対数曲線の回転体の体積 & \\
    \hline
  \end{tabular}
\end{table}

\begin{table}[htb]
  \centering
  \caption{解答}
  \renewcommand{\arraystretch}{1.4}
  \begin{tabular}{|c|c|p{9cm}|}
    \hline
    \textbf{問題番号} & \textbf{小問} & \textbf{答え} \\
    \hline
    第1問 & (1) & 斜辺の長さ$1-2r$ \\
    \cline{2-3}
    & (2) & 最大面積$\dfrac{1}{8}$ \\
    \hline
    第2問 & (1) & $f(a)$または$f(a+2^{n-1})$は$2^{n+1}$で割り切れる（証明問題） \\
    \cline{2-3}
    & (2) & 各$n$に対し条件を満たす$a_n$が存在（証明問題） \\
    \hline
    第3問 & -- & 対角線の交点は$AC$と$OB$の中点を結ぶ線分上にある（証明問題） \\
    \hline
    第4問 & -- & $\lim_{m\to\infty}\dfrac{L_m}{S_m}=1-\dfrac{1}{4a^2}$ \\
    \hline
    第5問 & (1) & $A(0),A(1),A(2),A(3)=42,36,0,6$ \\
    \cline{2-3}
    & (2) & 期待値$\dfrac{9}{14}$ \\
    \hline
    第6問 & (1) & $V(a)=\pi\left[u\{(\log u)^2-2\log u+2\}\right]_{u=1-a}^{3-a}$ \\
    \cline{2-3}
    & (2) & $a=2-\sqrt{2}$で最小値$2\pi\{(\log(1+\sqrt{2}))^2-2\sqrt{2}\log(1+\sqrt{2})+2\}$ \\
    \hline
  \end{tabular}
\end{table}

\paragraph{試験形態}
6問．試験時間は要確認．

\end{document}
